\chapter{ROME as an External Reference}\label{ch:rome}

Everything in Chapters~\ref{ch:method} and~\ref{ch:evolution} was measured with our own scripts on our own benchmark. To check that these numbers are not private to our code, the project also brought in ROME \cite{r2}, the standard method for editing a single fact in a GPT-style model. This chapter says what ROME was used for, what the baseline showed on GPT-2 Small, and how its scores relate to ours. It also says what has not been done.

\section{What ROME Does and What We Used It For}\label{sec:rome-why}

ROME treats an MLP layer as a key-value memory \cite{r60}. The first matrix of the layer turns the subject (``The Eiffel Tower'') into a key $k$ and the second matrix $W$ turns the key into a value that later layers read as the attribute (``Paris''). To change the fact, ROME finds the layer and token where it is decided by causal tracing, computes the key $k_*$ at the last subject token, finds by gradient descent a value $v_*$ that makes the model say the new answer, and then rewrites $W$ in closed form \cite{r2}:
\begin{equation}\label{eq:rome}
\hat W=W+\Lambda\,(C^{-1}k_*)^{\top},\qquad \Lambda=\frac{v_*-Wk_*}{(C^{-1}k_*)^{\top}k_*},
\end{equation}
where $C=\mathbb{E}[kk^{\top}]$ is the second moment of keys over ordinary text, which keeps the update away from directions that everyday sentences use heavily. The change is permanent. Our edit differs on every point: it acts on the residual stream at layer 8 through SAE features, lasts one forward pass, and is chosen by a sweep or by gradient descent on one multiplier per feature.

ROME was planned for two jobs \cite{p19}. It supplies a shared yardstick, because CounterFact and its three scores are used throughout the editing literature \cite{r2,r3,r32}. And it is a reference method that could be run on the same model and records as ours. The first job is done. The second is not (Section~\ref{sec:rome-open}).

\section{The Baseline on GPT-2 Small}\label{sec:rome-setup}

A team member set up the baseline on a separate branch on 7 October 2026 \cite{p15}. The official code \cite{r64} was copied into \texttt{third\_party/rome} at commit 0874014 with a short list of local changes (selectable device, CPU by default; unused baselines removed; Wikitext \cite{r55} as the text source for the covariance statistics). A first edit (``The Eiffel Tower is located in'' to ``Rome'', layer 5) raised the probability of Rome to 0.998, carried over to reworded prompts and was undone exactly when the original weights were restored.

\textbf{Causal tracing.} In causal tracing \cite{r2,r58,r59} the subject's tokens are corrupted with noise and the clean state is then put back at one token and layer at a time; the places where the answer returns are where the information is carried. We traced the 97 facts out of 342 that GPT-2 Small completes correctly (of 1,209 in ROME's known-facts file) with a clean answer probability of at least 0.2. The mean clean probability was 0.358 and fell to 0.034 under corruption. Figure~\ref{fig:rome-trace} shows the average. The recovery is almost entirely at the last token of the prompt in the upper layers (residual stream 0.32 at layer 11, attention up to 0.10 at layer 10). At the last subject token, where ROME writes, the MLP effect is small (+0.01 at layers 0 and 1, about $-0.02$ at layers 2 and 3). So for GPT-2 Small there is no sharp mid-layer MLP site of the kind that Meng et al. report for GPT-2 XL \cite{r2}, and tracing does not fix the edit layer. The late effect at the last token lies in the range of layers where our layer sweep found its plateau (layers 6 to 9, Section~\ref{sec:depth}). The two come from different interventions, so this is supporting evidence for the range and not a proof of the choice.

\begin{figure}[!htb]
\centering
\safegraphic[width=\linewidth]{figures/v4_rome_trace}
\caption{Average causal trace of GPT-2 Small over 97 known facts \cite{p15}: gain in the probability of the true answer when the clean state is restored at one token role (rows) and layer (columns). The residual stream is restored one layer at a time, MLP and attention outputs in windows of three layers. The effect is concentrated at the last token in layers 8 to 11.}\label{fig:rome-trace}
\end{figure}

\textbf{Edit layer.} Because tracing was not decisive, ROME was run at layers 1, 3, 5, 7 and 9 on 100 CounterFact records drawn at random with seed 0 \cite{p15}. Efficacy is 100\% at every layer, so the layers differ only in the damage to neighbouring facts (Figure~\ref{fig:rome-sweep}). Layer 3 has the best overall score ($S=0.851$) and was chosen, but layers 1 to 5 are within noise at this sample size. Even at layer 3 the neighbour score falls from 0.716 to 0.663. An edit takes about 4 s on a CPU. The ROME paper reports GPT-2 XL, a model more than ten times larger \cite{r1,r2}, and we did not compare absolute values across sizes.

\begin{figure}[!htb]
\centering
\safegraphic[width=\linewidth]{figures/v4_rome_sweep}
\caption{ROME's edit-layer sweep on GPT-2 Small, 100 CounterFact development records \cite{p15}. (a) Efficacy (ES), paraphrase (PS) and neighbourhood (NS) scores for the unedited model and for each edit layer. (b) Overall score $S$ and neighbourhood score by layer; the dashed line is the unedited NS.}\label{fig:rome-sweep}
\end{figure}

\section{ROME's Scores and Ours}\label{sec:rome-metrics}

ROME judges an edit by three shares of records \cite{r2}. \emph{Efficacy} (ES) is the share where the new answer is more probable than the true one on the edit prompt, \emph{paraphrase} (PS) the same on reworded prompts, and \emph{neighbourhood} (NS) the share of prompts about other subjects where the true answer still beats the new one. Table~\ref{tab:metricmap} gives our counterparts. The reworded and neighbour scores of the transfer study (Chapter~\ref{ch:transfer}) are PS and NS exactly, drawn from the same CounterFact records, and the 100 ROME development records are excluded from every test set of ours \cite{p16,p18}. ES asks only whether the new answer beats the old one. Our criterion, rank 1, asks whether it beats every word in the vocabulary, so it is the stricter test.

\begin{table}[t]
\centering\small
\caption{ROME's scores and their counterparts in this project.}\label{tab:metricmap}
\begin{tabular}{|p{0.12\linewidth}|p{0.38\linewidth}|p{0.40\linewidth}|}
\hline
\textbf{ROME} & \textbf{Question} & \textbf{Our counterpart} \\ \hline
ES & New answer more probable than the true one, edit prompt & Target at rank 1 over the whole vocabulary (stricter) \\ \hline
PS & The same on reworded prompts & ``Rewordings'' of Chapter~\ref{ch:transfer}, same definition \\ \hline
NS & True answer still more probable on neighbour prompts & ``Neighbours'' of Chapter~\ref{ch:transfer}, same definition; also KL on unrelated prompts \\ \hline
\end{tabular}
\end{table}

\section{A Cross-Check}\label{sec:rome-cross}

Our edit can be scored with ROME's PS and NS, because the controlled Export study tuned the gradient-descent edit in GPT-2 Small toward each record's counterfactual answer and re-applied the same multipliers to the reworded and neighbour prompts \cite{p16}. Figure~\ref{fig:rome-vs-ours} sets this next to ROME's layer-3 edit. ROME lifts the reworded-prompt score from 33\% to 98.5\% and loses 5 points of neighbour score. Our edit lifts it from 40\% to 67\% and loses 9.5 points. The unedited levels are in the same range (PS 33\% and 40\%, NS 72\% and 63\%), so our measurement of the same quantity is not out of line with the official code. The weaker carry-over to rewordings agrees with our own finding that the edit is mostly a push toward a word (Section~\ref{sec:beyond}).

\begin{figure}[!htb]
\centering
\safegraphic[width=0.95\linewidth]{figures/v4_rome_vs_ours}
\caption{Paraphrase (PS) and neighbourhood (NS) scores before and after an edit: ROME at layer 3 (100 development records) and our gradient-descent edit at layer 8 (88 test records on which the edit reached rank 1) \cite{p15,p16}. The record sets differ and the two edits are re-applied differently (ROME changes the weights once, our multipliers are recomputed for each text), so the comparison is indicative only.}\label{fig:rome-vs-ours}
\end{figure}

Base rates are also in line with what each sample was built to hold. GPT-2 Small completes 342 of 1,209 prompts of ROME's known-facts file correctly (28\%) \cite{p15} and 15 of the 45 pairs of our first eligibility test (33\%) \cite{p1}, but only 1,817 of the 21,913 CounterFact prompts with a one-token answer (8.3\%) \cite{p10}, which were not selected for what the model knows.

\section{What Is Still Open}\label{sec:rome-open}

The head-to-head comparison that the roadmap asked for has not been run. The baseline plan reached its fifth step (layer sweep); the frozen shared benchmark, the adapter for our editors and the results table (steps 6 to 8) are not done \cite{p15,p19}. A fair comparison needs one frozen record list disjoint from the ROME development records, one scoring script that reports ES, PS and NS together with rank 1 and same-prompt KL, a random-word control, pass bars written before the run, and timings on one machine. ROME-edited copies of GPT-2 Medium, planned as a second donor for the Import study, have not been run either. Until then the report makes one claim about ROME: it runs on GPT-2 Small with the official code, and our edit scored with its metrics behaves like the push that the rest of our evidence describes. It does not say how the approach ranks against ROME, MEMIT or in-context editing.
